% ============================================================
\section{Spectral Attribution Analysis}
\label{sec:analysis}
% ============================================================

Raw accuracy differences are insufficient to explain \emph{why} QIE
representations fail.
We attribute each failure mode to a spectral or geometric property of the
encoded feature matrix.

\subsection{Amplitude Encoding: Spectral Collapse on Low-Dimensional Data}
\label{sec:amplitude}

$\ell_2$-normalisation (Eq.~\ref{eq:amplitude}) projects all inputs onto the
unit hypersphere, destroying magnitude information.
The failure mode takes two related but distinct forms depending on the
intrinsic dimensionality of the input.

\paragraph{Rank collapse on tabular data.}
When the input features span a low-dimensional manifold — common in tabular
datasets with correlated features — normalisation concentrates the encoded
distribution on a small region of the unit sphere.
The encoded matrix collapses to near-rank-1.
Table~\ref{tab:spectral_amplitude} documents this for the tabular and
structured datasets.

\begin{table}[tbp]
\centering
\caption{Spectral properties of amplitude encoding.  On tabular datasets
failure is driven by rank collapse ($\erank \ll 1$); on image datasets rank
is adequate but conditioning is catastrophic.}
\label{tab:spectral_amplitude}
\small
\begin{tabular}{lrrl}
\toprule
Dataset & $\erank$ & $\log_{10}\kappa$ & Accuracy gap \\
\midrule
\multicolumn{4}{l}{\textit{Tabular / structured}} \\
Dry Bean        &     1.0 &  9.76 & $-0.67$ \\
Wine            &     1.4 &  3.82 & $-0.35$ \\
Breast Cancer   &     1.6 &  5.88 & $-0.20$ \\
Covertype       &     3.2 &  5.73 & $-0.27$ \\
\midrule
\multicolumn{4}{l}{\textit{Image (high-dimensional input)}} \\
CIFAR-10        &   460.9 &  4.48 & $-0.08$ \\
Fashion-MNIST   &   339.6 &  5.64 & $-0.04$ \\
\midrule
\multicolumn{4}{l}{\textit{Control}} \\
High-rank noise & 197.3   &  0.35 & $-0.00$ \\
\bottomrule
\end{tabular}
\end{table}

\paragraph{Ill-conditioning on image data.}
On high-dimensional image datasets, amplitude encoding does \emph{not} produce
rank collapse: Fashion-MNIST yields $\erank{=}339.6$ (out of $\dout{=}1{,}024$)
and CIFAR-10 yields $\erank{=}460.9$ (out of $\dout{=}4{,}096$).
Nevertheless, performance remains well below the best classical baseline.
The spectral explanation shifts from rank to conditioning: $\log_{10}\kappa{=}5.64$
for Fashion-MNIST and $4.48$ for CIFAR-10.
In the image case, $\ell_2$-normalisation imposes extreme scale variation
across the encoded dimensions, making the resulting Gram matrix ill-conditioned
even though its rank is not degenerate.
The failure mode is the same operation ($\ell_2$-normalisation) but acting on
intrinsically high-rank inputs; the symptom shifts from rank collapse to
catastrophic conditioning.

\paragraph{The control case.}
The high-rank noise dataset provides the falsifiable control: with $d{=}200$
isotropic Gaussian noise features, inputs already lie near-uniformly on the
unit sphere.
Normalisation neither collapses the rank ($\erank{=}197.3$) nor degrades
conditioning ($\kappa{=}2.22$), and amplitude is the only method achieving
near-parity with the best classical baseline ($d{=}{-}0.03$, $p{=}0.96$).

Figure~\ref{fig:kappa} shows the relationship between $\log_{10}\kappa$ and
accuracy gap across all datasets.

\begin{figure}[h!]
  \centering
  \includegraphics[width=0.75\textwidth]{figures/spectral_kappa_vs_gap.png}
  \caption{Condition number ($\log_{10}\kappa$) versus accuracy gap relative to
  best classical baseline for amplitude encoding across all ten datasets.
  The association holds regardless of whether failure is driven by rank collapse
  (low-$\erank$ tabular datasets) or ill-conditioning (high-$\erank$ image datasets).}
  \label{fig:kappa}
\end{figure}

\subsection{Angle Encoding: Geometric Equivalence to Linear}
\label{sec:angle}

The $(\cos(\theta/2), \sin(\theta/2))$ map is smooth and differentiable, but
under logistic regression it reduces to a scaled linear transform.
Table~\ref{tab:cka_angle} reports CKA between angle encoding and raw
(standardised) features.

\begin{table}[tbp]
\centering
\caption{CKA between angle encoding and raw linear features.}
\label{tab:cka_angle}
\small
\begin{tabular}{lr}
\toprule
Dataset & CKA(angle, raw linear) \\
\midrule
Fashion-MNIST   & 0.973 \\
High-rank noise & 0.985 \\
High-dim parity & 0.968 \\
Wine            & 0.971 \\
Breast Cancer   & 0.957 \\
Dry Bean        & 0.964 \\
CIFAR-10        & 0.971 \\
HIGGS           & 0.636 \\
Credit Fraud    & 0.514 \\
Covertype       & 0.527 \\
\bottomrule
\end{tabular}
\end{table}

On 7/10 datasets, CKA $\geq 0.95$: the angle-encoded Gram matrix is nearly
identical to that of standardised raw features.
The encoding is geometrically equivalent to a linear map under logistic
probing, providing no nonlinear representational lift.
The three lower-CKA datasets (HIGGS, Credit Fraud, Covertype) are exactly
those where magnitude or distributional structure breaks the
linear-scaling relationship; yet angle encoding does not improve performance
on these datasets either, consistent with the geometry being different rather
than better.

\subsection{Basis Encoding: Hamming Geometry Mismatch}
\label{sec:basis}

Basis encoding achieves variable $\erank$ across the benchmark: normalised
effective rank ($\erank/\dout$) ranges from 0.07 on CIFAR-10 to 0.88 on
high-dim parity, with tabular datasets concentrated in the 0.70--0.88 band.
Conditioning is favourable on image datasets ($\kappa{=}107$ on CIFAR-10,
$\kappa{=}142$ on Fashion-MNIST), where the 8-bit binary expansion preserves
distinct bit patterns across the high-dimensional feature space.
Its Gram matrix is geometrically distinct from all classical baselines.
On tabular datasets, $\mathrm{CKA}(\text{basis, raw linear})$ ranges from
0.36 to 0.45.
On image datasets the separation is more extreme: $\mathrm{CKA} = 0.011$ on
Fashion-MNIST and $0.001$ on CIFAR-10, confirming that 8-bit binary
quantisation produces a maximally orthogonal representation to pixel-space
linear features.
This geometric novelty, however, does not translate to better classification.
Binary quantisation imposes a Hamming geometry on the feature space: the
$\ell_2$ distance between two encoded samples reflects bit-flip counts, not
the smooth amplitude differences that discriminate classes on all ten
benchmark datasets.
As a result, basis encoding achieves 0/10 statistically competitive results
despite its favourable conditioning.

\subsection{The Spectral Prediction}
\label{sec:prediction}

Figure~\ref{fig:erank} plots effective rank (normalised by $\dout$) against
accuracy gap for all encodings and datasets.
The spectral attribution framework yields a single unified prediction: QIE
approaches competitive performance when $\erank / \dout \approx 1$ and
$\kappa \approx 1$ — that is, when the encoding does not collapse the input
geometry.
This condition is satisfied on one dataset (amplitude, high-rank noise) and
approached on none others, consistent with the prediction.

\begin{figure}[h!]
  \centering
  \includegraphics[width=0.75\textwidth]{figures/spectral_erank_vs_acc.png}
  \caption{Normalised effective rank ($\erank/\dout$) versus accuracy relative
  to raw linear baseline across all encodings and datasets.
  The one near-parity point (amplitude, high-rank noise) has
  $\erank/\dout{=}0.77$ and $\kappa{=}2.22$.}
  \label{fig:erank}
\end{figure}
